% Capítulo 5: Conclusiones
% Propósito: Sintetizar las conclusiones obtenidas en función de los objetivos alcanzados y plantear recomendaciones o trabajos futuros.
\section{Recomendaciones}

A partir del desarrollo de la presente investigación, se establecen recomendaciones orientadas al fortalecimiento de la plataforma web Comunica y al desarrollo de futuras investigaciones relacionadas con la gestión digital de incidencias y la optimización de los procesos de atención al usuario.

Se recomienda realizar un estudio de mercado dirigido entidades que ofertan servicios del municipio de Rivas, con el propósito de identificar entidades interesadas en adoptar Comunica. Este estudio permitirá conocer sus necesidades, evaluar la aceptación de la plataforma y determinar la viabilidad de su comercialización, considerando el precio propuesto y las condiciones de adquisición.

Se recomienda fortalecer las evaluaciones de usabilidad y accesibilidad mediante la participación de usuarios con diferentes niveles de conocimientos tecnológicos. Esto permitirá identificar dificultades de navegación, comprensión y utilización de las funcionalidades, contribuyendo a que Comunica ofrezca una experiencia de uso más intuitiva y accesible.

Asimismo, se recomienda realizar pruebas de rendimiento, seguridad y escalabilidad que permitan evaluar el comportamiento de la plataforma ante diferentes cargas de trabajo. Estas evaluaciones contribuirán a identificar posibles limitaciones técnicas y establecer las mejoras necesarias para garantizar la estabilidad, protección de la información y capacidad de crecimiento del sistema.

Para futuras investigaciones relacionadas con el desarrollo de plataformas de gestión de incidencias, se recomienda incorporar las buenas prácticas de ITIL 4, particularmente en los procesos de registro, clasificación, priorización, asignación, seguimiento y cierre de reportes. Su aplicación permitirá orientar la gestión de incidencias mediante procedimientos estructurados y criterios de evaluación que favorezcan la calidad y mejora continua del servicio.

Se recomienda que futuras investigaciones contemplen la evaluación del impacto operativo de las soluciones tecnológicas desarrolladas, mediante indicadores como los tiempos de respuesta, la cantidad de incidencias atendidas y la satisfacción de los usuarios. Esto permitirá obtener evidencia sobre la efectividad de estas herramientas y su contribución a la optimización de los procesos de atención.

De igual manera, se sugiere que otros investigadores adopten metodologías estructuradas como Design Science Research Methodology (DSRM), complementadas con marcos de desarrollo ágil como Scrum y prácticas de Extreme Programming (XP). Su aplicación permitirá mantener una relación coherente entre la identificación del problema, el diseño de la solución, su desarrollo y la evaluación de los resultados obtenidos.

Finalmente, se recomienda promover investigaciones relacionadas con la digitalización de los procesos de atención al usuario en instituciones públicas y empresas privadas de Nicaragua, considerando sus condiciones tecnológicas, organizacionales y sociales. Estos trabajos contribuirán a ampliar los antecedentes nacionales y generar conocimientos que puedan servir como referencia para el desarrollo de nuevas soluciones tecnológicas.
\newpage
\section{Concluciones}

La presente investigación permitió establecer las bases metodológicas, funcionales y tecnológicas para el desarrollo de la plataforma web Comunica, orientada a mejorar los procesos de registro, clasificación y seguimiento de incidencias. A través de la revisión bibliográfica y el estudio de mapeo sistemático, se identificaron soluciones tecnológicas existentes y sus principales aportes, lo que facilitó la definición de las funcionalidades necesarias para atender la problemática planteada.

En relación con el primer objetivo específico, se avanzó en el análisis y definición de los requerimientos funcionales y no funcionales mediante la revisión de investigaciones relacionadas con la gestión digital de incidencias y la identificación de las necesidades asociadas a los mecanismos tradicionales de atención. Este proceso permitió establecer los requisitos de autenticación, registro de reportes, geolocalización, clasificación, seguimiento y administración, así como elaborar un instrumento de recolección de datos para complementar el análisis de las necesidades de los usuarios.

Respecto al segundo objetivo específico, se definió una propuesta de arquitectura tecnológica basada en Angular, AdonisJS y PostgreSQL, junto con el diseño de los componentes necesarios para almacenar y gestionar la información de los reportes. Asimismo, se establecieron los elementos funcionales de las interfaces de usuario, considerando el registro de incidencias con evidencias fotográficas, ubicación geográfica y consulta de estados, proporcionando una estructura para la construcción e integración de la plataforma.

En cuanto al tercer objetivo específico, se estructuró el desarrollo de los módulos funcionales de Comunica mediante historias de usuario, matrices de trazabilidad y actividades organizadas en sprints. Esta planificación permitió relacionar los requerimientos establecidos con los componentes de registro, clasificación, seguimiento y administración de incidencias, incluyendo un panel de indicadores y visualizaciones geográficas. De esta manera, se definieron los criterios funcionales y técnicos que orientan la construcción y evaluación de la solución.

Finalmente, en relación con el objetivo general, la investigación permitió fundamentar y organizar el desarrollo de una plataforma web destinada a centralizar la recepción y gestión de incidencias en un entorno digital. La integración propuesta de mecanismos de registro, seguimiento y análisis de información constituye una base para mejorar la comunicación entre los usuarios y las entidades responsables. No obstante, la comprobación de su contribución efectiva a la optimización de los procesos de atención requiere completar las pruebas funcionales y realizar evaluaciones en condiciones reales de uso.